% =============================================================================
% INTRODUCTION TO VOLUME 6: HOLOGRAPHY, THE HLV BRIDGE AND THE STRUCTURE OF TIME
% (Vol. 6 of 8; continuation of Vol. 5)
% =============================================================================

\section*{Sixth Volume of the Document Collection}

With this volume the collection of individual documents on T0 theory / FFGFT
is continued beyond the original five volumes. Since the completion of
Volume~5 (Doc.~262, May 2026) the corpus has grown by more than 120
documents. They are included in Volumes 6, 7 and 8 in numerical order; the
boundary of each volume is placed so that a thematic break coincides with the
page limit of the print editions.

\subsection*{Content and Focus}

Volume~6 comprises Documents 263 to 310 from June and July 2026. The common
thread is the question of how the compact geometry $T^4$ passes over into the
observed three-dimensional, temporally open world, and what is preserved in
that passage in terms of information, phase and temporal structure. The
volume opens with the holographic principle in its fractal reading, the
cosmological degeneracy between $\Lambda$CDM and FFGFT, and the derivation of
the CMB peak ratios from the $T^4$. It continues with the projection chain
from $T^4$ to observable space, which leads through a series of comparison and
bridge documents on the Helix--Light--Vortex approach (HLV) up to the computed
3-versus-6 bridge in Doc.~285.

The third section deals with the dynamic sector: kinetic energy as a
$\xi$-fractal fraction of the static relation $\tilde T\cdot m=1$, the
separation of magnitude and phase in the four fermion sectors, and the twofold
origin of the angle $\theta=2/9$, dynamical (Doc.~291) and icosahedral
(Doc.~293). Documents 295 to 307 are devoted to time and information: the
time-winding as a logarithmic spiral, the reversibility test that identifies
information as an output rather than a base unit, the native time--energy
reciprocity, and the location of time within the state space. The volume
closes with the cosmic sector and the scale-anchor problem in comparison with
$\Lambda$CDM, and with resonance geometry as the overarching reference.

\subsection*{Organisation}

The 44 documents are divided into five parts: holography, cosmological
degeneracy and formalism correspondences (Docs.~263--269); change of
dimension, the HLV bridge and scale recursion (Docs.~270--285); dynamic
sector, magnitude and phase, $\theta=2/9$ (Docs.~286--294); time and
information (Docs.~295--307); cosmic sector and resonance geometry
(Docs.~308--310).

\subsection*{Reading the Documents}

As in the earlier volumes, each document is self-contained and is cited in
the text by its corpus number. The status markers [K], [B], [S] and
[SETZUNG] indicate the level of derivation reached; later corrections and
closed bridges are recorded in the corrections register Doc.~190, maintained
continuously in the repository (see the introduction to Volume~5). Where a later document refines or withdraws a statement of this
volume, the later version applies. Documents 308 to 310 were originally
typeset as multi-chapter treatises; in the book their chapters appear as
sections of the respective document.

\subsection*{Status}

The documents in this volume reflect their respective versions in the corpus
as of October 2026.
